\documentclass[10pt,a4paper]{amsart}
\usepackage[margin=19mm]{geometry}
\usepackage{amsmath,amssymb,mathtools}
\usepackage[scr=rsfs,cal=euler]{mathalfa}
\usepackage{xfrac}
\usepackage[unicode,hidelinks,bookmarks=false]{hyperref}
\newcommand{\ie}{i.e.\ }
\newcommand{\cf}{cf.\ }
\newcommand{\resp}{respectively\ }
\newcommand{\Subsection}{\subsection}
\providecommand{\nobreakdash}{\nobreak\hbox{-}}
\setlength{\emergencystretch}{2em}
\multlinegap0pt
\usepackage{amssymb}
\usepackage{mathtools}
\usepackage{algorithm}\usepackage{algpseudocode}
\newtheorem{remark}{Remark}
\newtheorem{theorem}{Theorem}
\newcommand{\wt}{\operatorname{wt}}
\title{Constraint-Preserving Genetic Algorithms for Embedding Linear Codes into Self-Orthogonal Codes}
\author{Haeun Lim, Junmin An, Jon-Lark Kim}
\date{Source: arXiv:2609.02135v1 (CC BY 4.0)}
\begin{document}
\maketitle

\noindent\emph{Source and licence.} Constraint-Preserving Genetic Algorithms for Embedding Linear Codes into Self-Orthogonal Codes. Version arXiv:2609.02135v1 (CC BY 4.0), distributed by arXiv under the Creative Commons Attribution 4.0 International licence, http://creativecommons.org/licenses/by/4.0/, as recorded in the arXiv licence metadata for this exact version and shown on the abstract page https://arxiv.org/abs/2609.02135v1. Full licence notice: \texttt{licenses/arxiv-2609.02135-CC-BY-4.0.txt}.\par\smallskip

\noindent\textit{Editorial scope.} Counted unit: the article's theorem orthogonal-group-generator (source0.tex lines 246--250). The article's complete original proof of it is delivered with it (source0.tex lines 251--273, one \texttt{proof} environment of 23 lines). No mathematics is written, repaired or re-derived: the statement and the proof are the article's own lines, in the article's own notation, and no retained line is substituted or re-flowed.  Retained uncounted as the source-local context the proof needs: lines 238--245.  Nothing was omitted from any retained span, so the package carries no omission record.  The retained lines cite 1 works, whose entries are transcribed verbatim from the article's own bibliography source in the e-print and delivered with short editorial labels recorded in the item's reference mapping.  No retained line contains a figure environment, an image-inclusion command or drawing code, so no image is delivered and the package declares no asset dependency.  Editorial formatting only, in addition: an AMS \texttt{amsart} preamble that restates the article's own declarations and macros used by the retained lines, namely the environments \texttt{remark}, \texttt{theorem} on separate counters, together with the standard packages those lines need.  The retained environments print in this document as Remark 1, Theorem 1 in order of appearance, whereas the article prints the same items with the numbering of its own full document.  The complete native source of the article as distributed in the arXiv e-print for this exact version is bundled unchanged as \texttt{sources/209176/source0.tex}.
\begin{remark}\label{orthogonal-group}
    It is known from \cite{Parametrization} that the orthogonal group $O(m,2)$ of order $m$ over $\mathbb{F}_2$ is given as follows:
    \begin{enumerate}
    \item if $1 \leq m \leq 3$, then $O(m,2) = \mathcal{P}_m$ where $\mathcal{P}_m$ is a permutation group of order $m$,
    \item if $m \geq 4$, then $O(m,2)=\langle\mathcal{P}_m, T_{\mathbf{u}}\rangle$
    where $\mathbf{u}$ is a vector of Hamming weight 4 over $\mathbb{F}_2$.
    \end{enumerate}
\end{remark}

\begin{theorem}\label{orthogonal-group-generator}
    For $m \geq 5$, $O(m,2)$ over $\mathbb{F}_2$ is generated by \[
    \mathcal{T} = \{T_\mathbf{u}~|~\wt(\mathbf{u})=4~for~\mathbf{u}\in\mathbb{F}_2^m\}.
    \]
\end{theorem}

\begin{proof}
    For any $\mathbf{u}_0 \in \mathbb{F}_2^m$, by Remark~\ref{orthogonal-group}, $O(m,2)=\langle\mathcal{P}_m, T_{\mathbf{u}_0}\rangle$. Thus, it is straightforward that $T_{\mathbf{u}_0} \in O(m,2)$.
    Next, we show that $\mathcal{T}$ generates $O(m,2)$, i.e., $\mathcal{P}_m \subseteq \langle\mathcal{T}\rangle$. Since $\mathcal{P}_m$ is generated by transpositions, it is sufficient to prove that any transposition $(ab)$ is generated by transvections in $\mathcal{T}$. We will show that any transposition can be written as a product of three transvections. Choose three distinct indices $c, d, e \notin \{a, b\}$, and set\[
    p=\mathbf{e}_a + \mathbf{e}_b + \mathbf{e}_d + \mathbf{e}_e,~~~~q=\mathbf{e}_b + \mathbf{e}_c + \mathbf{e}_d + \mathbf{e}_e.
    \] where $\mathbf{e}_i$ denotes the $i$-th standard basis vector of $\mathbb{F}_2^m$, that is, the vector whose $i$-th coordinate is $1$ and all other coordinates are $0$.
    Then $\wt(p)=\wt(q)=4$ and $p \cdot q = 1$. For $\mathbf{x} \in \mathbb{F}_2^m$, we have
\begin{align*}
    T_p T_q T_p(\mathbf{x}) &= T_p T_q(\mathbf{x} + (\mathbf{x} \cdot p)p)\\
    &= T_p (\mathbf{x} + (\mathbf{x} \cdot p)p + (\mathbf{x} \cdot q + \mathbf{x} \cdot p)q)\\
    &= \mathbf{x} + (\mathbf{x} \cdot p)p + (\mathbf{x} \cdot q + \mathbf{x} \cdot p)q + (\mathbf{x} \cdot q)p\\
    &= \mathbf{x} + (\mathbf{x} \cdot (p+q))(p+q)\\
    &= T_{p+q}(\mathbf{x})\\
    &= T_{\mathbf{e}_a+\mathbf{e}_b}(\mathbf{x}).
\end{align*}
For $\mathbf{x} \in \mathbb{F}_2^m$,
\begin{align*}
    T_{\mathbf{e}_a+\mathbf{e}_b}(\mathbf{x})&=\mathbf{x} + (\mathbf{x} \cdot (\mathbf{e}_a+\mathbf{e}_b))(\mathbf{e}_a+\mathbf{e}_b)\\
    &=\mathbf{x} + (x_a+x_b)(\mathbf{e}_a+\mathbf{e}_b)\\
    &=\mathbf{x} + (x_a+x_b)\mathbf{e}_a + (x_a+x_b)\mathbf{e}_b\\
\end{align*}
Hence, we conclude that $T_{\mathbf{e}_a+\mathbf{e}_b}$ interchanges the $a$-th and $b$-th coordinates of $\mathbf{x}$ and fixes all other coordinates.
Thus, $T_p T_q T_p = T_{\mathbf{e}_a+\mathbf{e}_b} = (ab)$.
\end{proof}

\begin{thebibliography}{99}
\bibitem[Parame]{Parametrization} G. J. Janusz, \emph{Parametrization of self-dual codes by orthogonal matrices}, Finite Fields Appl., 13(3) : 450–-491, 2007.
\end{thebibliography}
\end{document}
